% ======================================================================
% Appendix C — Self-contained verification core
% ======================================================================

\section{A self-contained verification core}\label{sec:self}

This appendix lets a reader with only this paper verify the sharpest
negative claims without the companion monograph. It collects the
elementary identity behind the four forms, the runner formulation with
its certificate, the explicit statement of the sum-only equality, and
one-line hand checks of the headline witnesses.

\subsection{The identity and the runner formulation}

\begin{lemma}[depth identity]\label{lem:depthid}
For every $t\in\T$,
$\max_{1\le j\le n}\wnc{\tfrac12-v_jt}=\tfrac12-\min_{1\le j\le n}\wnc{v_jt}$.
Consequently $\delta(V)=\tfrac12-\lambda(V)$ with
$\delta(V)=\operatorname{dist}_\infty\!\bigl(\R v,\;\Z^n+\tfrac12\mathbf{1}
\bigr)$, and the $\delta$-form of \S\ref{sec:prelim} is exactly
$\lrc_n$.
\end{lemma}

\begin{proof}
Each map $x\mapsto\wnc{x}$ satisfies
$\wnc{\tfrac12-x}=\tfrac12-\wnc{x}$: reflecting about $\tfrac12$
exchanges the two half-intervals and fixes the midpoint. Taking $x=v_jt$
and maximizing over $j$ gives the displayed identity (a maximum of
quantities each equal to $\tfrac12$ minus the corresponding term is
$\tfrac12$ minus the minimum). The distance form is the same statement
along the trajectory line $tv$, $t\in\R$.
\end{proof}

The computations of this paper use the \emph{runner formulation}
\eqref{eq:runnerform}, restated here:
\[
D(c)\;=\;\min_{\sigma\in[0,1)}\;\max\Bigl(\wnc{\sigma},\;
\max_{2\le j\le n}\wnc{c_j-v_j\sigma}\Bigr).
\]
Two facts make it auditable. First, it is the quotient-norm depth
$\operatorname{dist}_{\mathcal{X}_v}(c,\Lambda)$ written in the
$c$-coordinates of \S\ref{sec:prelim}: the linear forms $\phi^{1j}$,
$\phi^{ij}$ listed there generate the norm, the parameter $\sigma$ is the
free coordinate $x_1$ of a lift, and minimizing over $\sigma$ minimizes
over the lift. Second, it is \emph{Lipschitz-auditable}: each term
$\wnc{c_j-v_j\sigma}$ is $v_j$-Lipschitz in $\sigma$ and the whole
objective is piecewise linear with breakpoints at the half-integers
$c_j-\tfrac12-v_j\sigma\in\Z$, so the minimum is attained at a rational
breakpoint or crossing and is exactly computable by hand or by exact
rational arithmetic. The independent linear-form/lattice-search
implementation of the program agrees with it exactly on every instance
used in this paper (forty random rational points per instance).

\subsection{The box certificate (Lipschitz in \texorpdfstring{$c$}{c})}

\begin{lemma}[sup-norm Lipschitz]\label{lem:clip}
For all $c,c'\in\T^{n-1}$,
$\lvert D(c)-D(c')\rvert\le\lVert c-c'\rVert_\infty$.
\end{lemma}

\begin{proof}
Fix $\sigma$. Each term $\wnc{c_j-v_j\sigma}$ differs from
$\wnc{c'_j-v_j\sigma}$ by at most $|c_j-c'_j|$ (the torus norm is
$1$-Lipschitz), the $\wnc\sigma$ term is unchanged, and the maximum of
finitely many functions each moved by at most
$\lVert c-c'\rVert_\infty$ is itself moved by at most
$\lVert c-c'\rVert_\infty$. Hence
$G(c,\sigma)\le G(c',\sigma)+\lVert c-c'\rVert_\infty$ for every
$\sigma$, and minimizing over $\sigma$ preserves the inequality in
both directions.
\end{proof}

\begin{corollary}[box certificate]\label{cor:box}
Let $C$ be a box with center $c_0$ and half-width
$h=\max_j\tfrac12(\overline{c_j}-\underline{c_j})$. Then
$\max_{c\in C}D(c)\le D(c_0)+h$. In particular, if $D(c_0)+h\le r$
with $D(c_0)$ computed exactly, then $C$ contains no point of depth
above $r$.
\end{corollary}

This is the test the first $m=10$ attempt lacked. At the residual
width it left behind---boxes of $1/512$ per side, so $h=1/1024$, with
exact center depths at most $1135/3584$---the corollary gives
\[
\max_{c\in C}D(c)\;\le\;\tfrac{1135}{3584}+\tfrac{1}{1024}
\;=\;\tfrac{2277}{7168}\;\approx\;0.31766\;<\;\tfrac{7}{22},
\]
closing the gap that the minimax bound $U(C)$ could not. The v3
branch-and-bound (\S\ref{sec:bb}) interleaves exactly this test with
the exact minimax
test; the adversarial soundness check is worth recording here because
it is a one-line computation in the same style as the hand checks
below: the $1/512$-box around the $6/19$ witness at
$v=(1,2,3,4,15)$, $r=5/16$, has exact center depth $6133/19456$, and
$6133/19456+1/1024>5/16$---the certificate correctly refuses a box
that contains a true deeper hole. The corresponding checks at the
harmonic witnesses ($347/1024$ and $1045/3072$ center depths against
$r=1/3$) refuse identically.

\subsection{The sum-only equality, explicitly}

\begin{theorem}[sum-only equality; explicit form]\label{thm:equality-explicit}
Let $V=(v_1,\dots,v_n)$ have distinct positive entries with $\gcd(V)=1$
and $0<M(V)<\tfrac12$. For each unordered pair $\{u,w\}\subset[n]$ write
$N_{uw}=v_u+v_w$ and
\[
\tau(u,w)\;=\;\max_{0\le k<N_{uw}}\;\min_{p\in[n]\setminus\{w\}}\;
\wnc{\frac{v_pk}{N_{uw}}}.
\]
Then $M(V)=\max_{\{u,w\}}\tau(u,w)$, and consequently $\lrc_n$ holds
for all $V$ if and only if for every $V$ there is a pair $\{u,w\}$ with
$\tau(u,w)\ge 1/(n+1)$.
\end{theorem}

The two ingredients (proved in the companion monograph, Theorem~1.1
there; the finitization is what every instance of \S\ref{sec:tm-covering}
entered through) are: (i) at a global optimum $t^\ast$ with
$0<M(V)<\tfrac12$ the binding runners can be chosen two-sided---one
rising, one falling---which forces $(v_u+v_w)t^\ast\in\Z$ for that pair,
so $t^\ast$ lies on the \emph{sum} lattice of denominator $v_u+v_w$; and
(ii) on that lattice the two binding runners are exact reflections
($\wnc{v_uk/N}=\wnc{v_wk/N}$ for all $k$, since $v_u\equiv-v_w \pmod N$),
so the $n$-runner minimum equals an $(n-1)$-runner minimum. The
companion's Lemma~2.1 also records the boundary cases $M(V)=\tfrac12$
(all speeds odd) and $M(V)=0$ (impossible).

\subsection{Hand checks of the headline witnesses}

All three strict harmonic-family witnesses are one-minute hand
computations with \eqref{eq:runnerform}.

\emph{The $97/288$ witness.} Take $v=(1,2,3,4,5)$ and
$c=(\tfrac1{32},\tfrac{15}{32},\tfrac3{32},\tfrac78)$. At
$\sigma=\tfrac{31}{288}$ (in $288$-ths, with $c=(\tfrac9{288},
\tfrac{135}{288},\tfrac{27}{288},\tfrac{252}{288})$):
$\wnc\sigma=\tfrac{31}{288}$;
$\wnc{c_2-2\sigma}=\wnc{\tfrac{9-62}{288}}=\tfrac{53}{288}$;
$\wnc{c_3-3\sigma}=\wnc{\tfrac{135-93}{288}}=\tfrac{42}{288}$;
$\wnc{c_4-4\sigma}=\wnc{\tfrac{27-124}{288}}=\tfrac{97}{288}$;
$\wnc{c_5-5\sigma}=\wnc{\tfrac{252-155}{288}}=\tfrac{97}{288}$.
The maximum is $\tfrac{97}{288}$, and it is minimized: runners $4$ and
$5$ have opposite slopes ($+4$, $-5$), so any local shift of $\sigma$
increases one of them. Hence $D(c)=97/288>\tfrac13=\mathrm{thr}_5$.

\emph{The $607/1792$ witness.} Take
$c=(\tfrac{11}{256},\tfrac{163}{256},\tfrac{15}{256},
\tfrac{79}{256})$ and $\sigma=\tfrac{89}{896}$: the five terms are
$\tfrac{89}{896}$, $\tfrac{279}{1792}$, $\tfrac{607}{1792}$,
$\tfrac{607}{1792}$, $\tfrac{337}{1792}$, so $D(c)=607/1792>
\tfrac13$, the branch-and-bound follow-up's improvement (two terms tie
at the maximum with opposite effective slopes).

\emph{The audit's $4183/12288$ witness.} Take
$c=(\tfrac{465}{4096},\tfrac{605}{1024},\tfrac{1019}{4096},
\tfrac{315}{4096})$ and $\sigma=\tfrac{905}{6144}$: the five terms are
\[
\tfrac{905}{6144},\quad \tfrac{2225}{12288},\quad
\tfrac{1830}{12288},\quad \tfrac{4183}{12288},\quad
\tfrac{4183}{12288},
\]
so $D(c)=4183/12288>\tfrac13$. This point was found during the independent audit of this
paper (forty random local searches at the harmonic instance, exact
re-evaluation), and we adopt it as the harmonic family's recorded lower
bound; its discovery is itself a datum for \S\ref{sec:bb}'s discussion
of search thoroughness.

The remaining strict witnesses ($29/80$ at $v=(1,2,3,4,5,6)$ and $4/11$
at the non-extremal $v=(1,2,3,4,5,7)$) check identically: substitute the
$c$-coordinates of Table~\ref{tab:refute}, evaluate the five (resp.\ \
six) terms at the recorded $\sigma$, and read off the balanced maximum.

\emph{The $7/22$ torsion center (the $m=10$ certificate's witness).}
Take $v=(1,2,3,4,10)$ and $c=(\tfrac12,0,\tfrac12,\tfrac12)$. At
$\sigma=\tfrac{7}{22}$ the five terms of \eqref{eq:runnerform} are
$\tfrac{7}{22}$, $\tfrac{3}{22}$, $\tfrac{1}{22}$, $\tfrac{5}{22}$,
$\tfrac{7}{22}$: explicitly $\wnc\sigma=\tfrac7{22}$;
$\wnc{\tfrac12-2\sigma}=\tfrac3{22}$;
$\wnc{-3\sigma}=\tfrac1{22}$;
$\wnc{\tfrac12-4\sigma}=\tfrac5{22}$; and
$\tfrac12-10\sigma=-\tfrac{59}{22}\equiv\tfrac7{22}\pmod 1$, so the
fifth term is $\tfrac7{22}$. The maximum is $\tfrac7{22}$, and it is
minimized: the binding terms are the $\sigma$-term (slope $+1$) and
the fifth term (slope $-10$), opposite signs, so any local shift of
$\sigma$ increases one of them. Hence $D(c^\ast)=7/22$, and the v3
certificate of \S\ref{sec:bb} ($\rho\le7/22$) upgrades this to
$\rho(1,2,3,4,10)=\tfrac{7}{22}$ exactly.

Every value in Table~\ref{tab:refute} was additionally re-asserted by
two independent implementations in the provenance run
(Appendix~\ref{sec:prov}).
